\begin{ouverture}[courbe]{4}{Équations et inéquations}{Partie 1 · Nombres et calculs}
{Résoudre · Modéliser · Étudier un signe}
\objectif{Donner l’ensemble exact des solutions.}
\objectif{Utiliser produit nul et tableaux de signes.}
\objectif{Interpréter les solutions dans une situation concrète.}
\prerequis{Calcul littéral et intervalles.}
\end{ouverture}
\pagestyle{otMain}\essentiel\fontsize{11}{15}\selectfont
\essSection{1}{Équations : transformer sans perdre de solution}
\begin{definition}Résoudre une équation, c’est trouver toutes les valeurs autorisées qui rendent l’égalité vraie. Une transformation équivalente conserve exactement cet ensemble.\end{definition}
\begin{propriete}On peut ajouter un même nombre aux deux membres, ou les multiplier par un même nombre non nul. Un produit est nul si et seulement si au moins l’un de ses facteurs est nul.\end{propriete}
\begin{exemple}$(2x-3)(x+4)=0$ donne $2x-3=0$ ou $x+4=0$, donc $S=\{-4;3/2\}$.\end{exemple}
\begin{propriete}$x^2=a$ n’a pas de solution réelle si $a<0$ ; sa seule solution est 0 si $a=0$ ; pour $a>0$, les solutions sont $-\sqrt a$ et $\sqrt a$.\end{propriete}
\begin{savoirfaire}[Équation quotient]Noter d’abord les valeurs interdites. Pour $(x+1)/(x-2)=3$, imposer $x\ne2$, puis écrire $x+1=3(x-2)$. On trouve $x=7/2$, autorisé.\end{savoirfaire}
\begin{attention}Diviser par une expression contenant $x$ peut faire perdre les valeurs pour lesquelles elle est nulle. Préférer la factorisation.\end{attention}
\clearpage
\essSection{2}{Inégalités et signes}
\begin{propriete}Additionner une même quantité conserve l’ordre. Multiplier ou diviser par un nombre strictement négatif renverse le sens de l’inégalité. On peut additionner deux inégalités de même sens.\end{propriete}
\begin{exemple}$-3x+2\ge8$ donne $-3x\ge6$, puis $x\le-2$.\end{exemple}
\begin{propriete}[Comparer deux nombres]Comparaison additive : $a\le b$ équivaut à $a-b\le0$. Comparaison multiplicative, pour $a>0$ et $b>0$ : $a\le b$ équivaut à $a/b\le1$. La différence $b-a$ s’interprète comme une variation additive, le rapport $b/a$ comme un coefficient multiplicateur.\end{propriete}
\begin{savoirfaire}[Produit ou quotient]
Déterminer les zéros de chaque facteur et les valeurs interdites. Les placer dans l’ordre dans un tableau. Sur chaque intervalle, appliquer la règle des signes. Inclure les zéros si l’inégalité est large, mais exclure toujours un dénominateur nul.
\end{savoirfaire}
\begin{exemple}
Pour $(x-1)(x+2)\le0$, les changements de signe se font en $-2$ et 1. Le produit est négatif entre ces valeurs : les solutions sont $[-2;1]$.
\end{exemple}
\begin{savoirfaire}[Modéliser]
Choisir la variable et son domaine ; traduire les contraintes ; résoudre ; revenir à la question. Pour un nombre de billets, ne garder que les entiers autorisés. Une borne réelle n’est pas toujours une réponse concrète.
\end{savoirfaire}
\clearpage\rubrique{Découverte}
\begin{blocLivre}{ex-4-1}
\begin{exercice}[1]{Une équation linéaire}
Résoudre $3x-7=11$ et vérifier la solution.
\end{exercice}
\end{blocLivre}
\begin{blocLivre}{ex-4-2}
\begin{exercice}[1]{L’inconnue des deux côtés}
Résoudre $5(x-2)=2x+8$.
\end{exercice}
\end{blocLivre}
\begin{blocLivre}{ex-4-3}
\begin{exercice}[1]{Un signe à renverser}
Résoudre $-4x+3>15$ et donner l’intervalle solution.
\end{exercice}
\end{blocLivre}
\begin{blocLivre}{ex-4-4}
\begin{exercice}[1]{Un carré donné}
Résoudre dans $\mathbb R$ les trois équations $x^2=49$, $x^2=0$ et $x^2=-9$.
\end{exercice}
\end{blocLivre}
\begin{blocLivre}{ex-4-5}
\begin{exercice}[1]{Produit nul}
Résoudre $(x-5)(2x+3)=0$.
\end{exercice}
\end{blocLivre}

\clearpage\rubrique{Entraînement}
\begin{blocLivre}{ex-4-6}
\begin{exercice}[2]{Une équation quotient}
Résoudre $\dfrac{2x+1}{x-3}=5$.
\end{exercice}
\end{blocLivre}
\begin{blocLivre}{ex-4-7}
\begin{exercice}[2]{Tout, rien ou une valeur}
Résoudre séparément $2(x+1)=2x+2$ et $2(x+1)=2x+3$.
\end{exercice}
\end{blocLivre}
\begin{blocLivre}{ex-4-8}
\begin{exercice}[2]{Un produit de deux facteurs}
Étudier le signe de $(x-1)(x+3)$ puis résoudre $(x-1)(x+3)\le0$.
\label{fig-ex-4-8}
\figureGeo{signes-produit}
\end{exercice}
\end{blocLivre}
\begin{blocLivre}{ex-4-9}
\begin{exercice}[2]{Un quotient et sa valeur interdite}
Résoudre $\dfrac{x-2}{x+1}>0$.
\label{fig-ex-4-9}
\figureGeo{signes-quotient}
\end{exercice}
\end{blocLivre}
\begin{blocLivre}{ex-4-10}
\begin{exercice}[2]{Déplacer une expression}
Résoudre $(x-2)^2\le(x-2)(x+4)$.
\end{exercice}
\end{blocLivre}

\clearpage\rubrique{Problèmes}
\begin{blocLivre}{ex-4-11}
\begin{exercice}[3]{Un abonnement avantageux}
Une salle propose 8 euros par séance, ou un abonnement à 30 euros puis 5 euros par séance. À partir de combien de séances l’abonnement est-il strictement moins cher ?
\end{exercice}
\end{blocLivre}
\begin{blocLivre}{ex-4-12}
\begin{exercice}[3]{Un budget avec remise fixe}
Une commande coûte 4,50 euros par carnet et 7 euros de livraison. On dispose de 52 euros. Combien de carnets au maximum peut-on commander ?
\end{exercice}
\end{blocLivre}
\begin{blocLivre}{ex-4-13}
\begin{exercice}[3]{Des dimensions à retrouver}
Un rectangle a pour longueur $x+3$ et largeur $x-1$, en centimètres, avec $x>1$. Son aire vaut 21. Résoudre sans utiliser de discriminant.
\label{fig-ex-4-13}
\figureGeo{rectangle-equation}
\end{exercice}
\end{blocLivre}
\begin{blocLivre}{ex-4-14}
\begin{exercice}[3]{Un paramètre au premier degré}
Pour un réel $m$, discuter le nombre de solutions de $(m-2)x=6$. Donner la solution lorsqu’elle existe.
\end{exercice}
\end{blocLivre}
\begin{blocLivre}{ex-4-15}
\begin{exercice}[3]{Une équation qui perdrait zéro}
Résoudre $x(x-3)=2x$. Expliquer pourquoi diviser immédiatement par $x$ conduit à une réponse incomplète.
\end{exercice}
\end{blocLivre}

\clearpage
